\documentclass[aspectratio=169]{beamer}

% ── Tema ──────────────────────────────────────────────────────────────────────
\usetheme{Madrid}
\usecolortheme{seahorse}
\setbeamertemplate{navigation symbols}{}
\setbeamertemplate{footline}[frame number]

% ── Pacotes ───────────────────────────────────────────────────────────────────
\usepackage[utf8]{inputenc}
\usepackage[T1]{fontenc}
\usepackage[brazil]{babel}
\usepackage{amsmath, amssymb}
\usepackage{listings}
\usepackage{xcolor}
\usepackage{tikz}
\usetikzlibrary{arrows.meta, positioning, shapes.geometric}
\usepackage{booktabs}
\usepackage{multirow}
\usepackage{graphicx}

% ── TikZ styles (definidos no preambulo para evitar conflito com beamer #1) ───
\tikzset{
  pipeblue/.style={draw, rounded corners, fill=blue!20, minimum height=0.75cm,
                   minimum width=3cm, text centered, font=\small, inner sep=3pt},
  pipeorange/.style={draw, rounded corners, fill=orange!20, minimum height=0.75cm,
                     minimum width=3cm, text centered, font=\small, inner sep=3pt},
  pipeorangedash/.style={draw, dashed, rounded corners, fill=orange!10, minimum height=0.75cm,
                         minimum width=3cm, text centered, font=\small, inner sep=3pt},
  pipepurple/.style={draw, rounded corners, fill=purple!15, minimum height=0.75cm,
                     minimum width=3cm, text centered, font=\small, inner sep=3pt},
  pipeyellow/.style={draw, rounded corners, fill=yellow!30, minimum height=0.75cm,
                     minimum width=3cm, text centered, font=\small, inner sep=3pt},
  pipered/.style={draw, rounded corners, fill=red!10, minimum height=0.75cm,
                  minimum width=3cm, text centered, font=\small, inner sep=3pt},
  pipegreen/.style={draw, rounded corners, fill=green!20, minimum height=0.75cm,
                    minimum width=3cm, text centered, font=\small, inner sep=3pt},
  pipegreendk/.style={draw, rounded corners, fill=green!30, minimum height=0.75cm,
                      minimum width=3cm, text centered, font=\small, inner sep=3pt},
  motblue/.style={draw, rounded corners, fill=blue!15, minimum width=2.8cm,
                  minimum height=0.6cm, text centered, font=\small},
  motorange/.style={draw, rounded corners, fill=orange!20, minimum width=2.8cm,
                    minimum height=0.6cm, text centered, font=\small},
  motred/.style={draw, rounded corners, fill=red!10, minimum width=2.8cm,
                 minimum height=0.6cm, text centered, font=\small},
  motgreen/.style={draw, rounded corners, fill=green!20, minimum width=2.8cm,
                   minimum height=0.6cm, text centered, font=\small},
  motgreendk/.style={draw, rounded corners, fill=green!30, minimum width=2.8cm,
                     minimum height=0.6cm, text centered, font=\small},
  piparr/.style={-Latex, thick, esbmcblue},
  piplbl/.style={font=\scriptsize, text=gray},
}

% ── Cores ─────────────────────────────────────────────────────────────────────
\definecolor{esbmcblue}{RGB}{0,82,156}
\definecolor{codegray}{RGB}{245,245,245}
\definecolor{codegreen}{RGB}{0,128,0}
\definecolor{codepurple}{RGB}{128,0,128}
\definecolor{okgreen}{RGB}{0,140,0}
\definecolor{failred}{RGB}{180,0,0}

% ── Listings ──────────────────────────────────────────────────────────────────
\lstdefinestyle{cstyle}{
  language=C,
  basicstyle=\ttfamily\scriptsize,
  backgroundcolor=\color{codegray},
  keywordstyle=\color{esbmcblue}\bfseries,
  commentstyle=\color{codegreen},
  stringstyle=\color{codepurple},
  frame=single,
  rulecolor=\color{gray!40},
  breaklines=true,
  tabsize=2,
  showstringspaces=false,
}
\lstdefinestyle{pystyle}{
  language=Python,
  basicstyle=\ttfamily\scriptsize,
  backgroundcolor=\color{codegray},
  keywordstyle=\color{esbmcblue}\bfseries,
  commentstyle=\color{codegreen},
  stringstyle=\color{codepurple},
  frame=single,
  rulecolor=\color{gray!40},
  breaklines=true,
  tabsize=2,
  showstringspaces=false,
}

% ── Metadados ─────────────────────────────────────────────────────────────────
\title[Verificação Formal de Redes Neurais]{%
  Verificação Formal de Redes Neurais com ESBMC}

\subtitle{Pipeline: PyTorch $\to$ ONNX $\to$ Quantização $\to$ Verificação SMT}

\author{Matheus Serrão Uchoa}

\institute{Laboratório de Verificação Formal --- PIBIC}

\date{2026}

% ══════════════════════════════════════════════════════════════════════════════
\begin{document}
% ══════════════════════════════════════════════════════════════════════════════

\begin{frame}
  \titlepage
\end{frame}

\begin{frame}{Sumário}
  \tableofcontents
\end{frame}

% ══════════════════════════════════════════════════════════════════════════════
\section{Motivação e Contexto}
% ══════════════════════════════════════════════════════════════════════════════

\begin{frame}[fragile]{Por que verificar formalmente redes neurais?}
  \begin{columns}
    \column{0.55\textwidth}
    \textbf{Problema}: redes neurais em sistemas críticos
    \begin{itemize}
      \item Veículos autônomos
      \item Diagnóstico médico
      \item Sistemas embarcados (ESP32, MCUs)
    \end{itemize}
    \vspace{0.4em}
    \textbf{Limitação do teste clássico}:
    \begin{itemize}
      \item Testa \emph{amostras} do espaço de entrada
      \item Não garante comportamento para \emph{todas} as entradas
    \end{itemize}
    \vspace{0.4em}
    \begin{alertblock}{Verificação formal com ESBMC}
      \centering Prova matemática de propriedades\\para \textbf{todo} input válido
    \end{alertblock}

    \column{0.42\textwidth}
    \begin{tikzpicture}[scale=0.9]
      \node[motblue]    (train)  at (0,3.5) {Treinamento};
      \node[motorange]  (test)   at (0,2.5) {Teste (amostras)};
      \node[motred]     (deploy) at (0,1.5) {\textcolor{failred}{Deploy (risco)}};
      \node[motgreen]   (verify) at (0,0.5) {\textcolor{okgreen}{Verif.\ Formal}};
      \node[motgreendk] (safe)   at (0,-0.3) {\textcolor{okgreen}{Deploy seguro}};
      \draw[-Latex] (train)  -- (test);
      \draw[-Latex] (test)   -- (deploy);
      \draw[-Latex, dashed, esbmcblue, thick] (test) -- ++(1.8,0) |- (verify);
      \draw[-Latex] (verify) -- (safe);
    \end{tikzpicture}
  \end{columns}
\end{frame}

% ──────────────────────────────────────────────────────────────────────────────
\begin{frame}{ESBMC --- Bounded Model Checker}
  \begin{columns}
    \column{0.5\textwidth}
    \textbf{ESBMC} --- Efficient SMT-Based BMC:
    \begin{itemize}
      \item Verifica programas C/C++ via SMT
      \item Solvers: Z3, Boolector, Bitwuzla
      \item Primitivas essenciais:
        \begin{itemize}
          \item \texttt{\_\_ESBMC\_assume(cond)} --- restringe busca
          \item \texttt{\_\_ESBMC\_assert(cond, msg)} --- propriedade
        \end{itemize}
    \end{itemize}
    \vspace{0.5em}
    \begin{block}{Por que ponto fixo em vez de float?}
      \texttt{--floatbv} usa teoria IEEE 754 --- exponencialmente mais difícil.
      Ponto fixo reduz a \textbf{bit-vector arithmetic}: Z3/Boolector
      resolvem em segundos.
    \end{block}

    \column{0.47\textwidth}
    \begin{exampleblock}{Fluxo de verificação}
      \begin{enumerate}
        \item Pesos quantizados $\to$ constantes C
        \item \texttt{nondet\_int()} para entradas simbólicas
        \item \texttt{\_\_ESBMC\_assume} para domínio
        \item \texttt{\_\_ESBMC\_assert} para propriedades
        \item Solver BV prova (ou refuta) formalmente
      \end{enumerate}
    \end{exampleblock}
  \end{columns}
\end{frame}

% ══════════════════════════════════════════════════════════════════════════════
\section{Caso de Estudo 1 --- MLP XOR}
% ══════════════════════════════════════════════════════════════════════════════

\begin{frame}[fragile]{MLP XOR --- Arquitetura e Problema}
  \begin{columns}
    \column{0.5\textwidth}
    \textbf{Rede}: MLP 2\;$\to$\;4\;$\to$\;1
    \begin{itemize}
      \item Tarefa: classificar XOR binário
      \item Oculta: ReLU \quad Saída: Sigmoid
      \item Framework: PyTorch
    \end{itemize}
    \vspace{0.5em}
    \begin{center}
      \begin{tabular}{cc|c}
        \toprule
        $x_1$ & $x_2$ & XOR \\
        \midrule
        0 & 0 & \textcolor{failred}{0} \\
        0 & 1 & \textcolor{okgreen}{1} \\
        1 & 0 & \textcolor{okgreen}{1} \\
        1 & 1 & \textcolor{failred}{0} \\
        \bottomrule
      \end{tabular}
    \end{center}

    \column{0.47\textwidth}
    \begin{tikzpicture}[scale=0.85,
      nd/.style={circle, draw, fill=blue!15, minimum size=0.65cm, font=\scriptsize}]
      \node[nd] (i1) at (0,2.25) {$x_1$};
      \node[nd] (i2) at (0,0.75) {$x_2$};
      \node[nd] (h1) at (2,3.5)  {$h_1$};
      \node[nd] (h2) at (2,2.25) {$h_2$};
      \node[nd] (h3) at (2,1.0)  {$h_3$};
      \node[nd] (h4) at (2,-0.25){$h_4$};
      \node[nd, fill=green!20] (o) at (4,1.5) {$\hat{y}$};
      \foreach \src in {i1,i2}
        \foreach \dst in {h1,h2,h3,h4}
          \draw[gray!60, ->] (\src) -- (\dst);
      \foreach \src in {h1,h2,h3,h4}
          \draw[gray!60, ->] (\src) -- (o);
      \node[above=0.1cm] at (0,2.9)  {\scriptsize Input};
      \node[above=0.1cm] at (2,3.9)  {\scriptsize Hidden};
      \node[above=0.1cm] at (4,2.1)  {\scriptsize Output};
    \end{tikzpicture}
  \end{columns}
\end{frame}

% ──────────────────────────────────────────────────────────────────────────────
\begin{frame}[fragile]{Pipeline Completo --- MLP XOR}
  \begin{tikzpicture}
    \node[pipeblue]       (tr) at (0,0)    {\texttt{mlp\_training.py}};
    \node[pipeorange]     (on) at (4.2,0)  {\texttt{mlp\_model.onnx}};
    \node[pipepurple]     (ex) at (8.4,0)  {\texttt{onnx\_mlp\_extractor}};
    \node[pipeyellow]     (qp) at (8.4,-2) {\texttt{quantize\_mlp.py}};
    \node[pipered]        (hc) at (8.4,-4) {\texttt{verify\_mlp\_qnn.c}};
    \node[pipegreen]      (es) at (4.2,-4) {\textbf{ESBMC} / Boolector};
    \node[pipegreendk]    (ok) at (0,-4)   {\textcolor{okgreen}{\textbf{SUCCESSFUL}}};
    \node[pipeorangedash] (wh) at (4.2,-2) {\texttt{mlp\_weights.h}};

    \draw[piparr] (tr) -- (on)  node[above,piplbl,midway]{export};
    \draw[piparr] (on) -- (ex)  node[above,piplbl,midway]{parse};
    \draw[piparr, dashed] (tr)  -- (wh) node[left,piplbl,midway]{\scriptsize alternativa};
    \draw[piparr] (ex) -- (qp)  node[right,piplbl,midway]{float};
    \draw[piparr, dashed] (wh)  -- (qp) node[below,piplbl,midway]{\scriptsize{-{}-source h}};
    \draw[piparr] (qp) -- (hc)  node[right,piplbl,midway]{Q8.8};
    \draw[piparr] (hc) -- (es);
    \draw[piparr] (es) -- (ok);
  \end{tikzpicture}
\end{frame}

% ──────────────────────────────────────────────────────────────────────────────
\begin{frame}[fragile]{Etapa 1 --- Treinamento PyTorch}
  \begin{columns}
    \column{0.48\textwidth}
    \textbf{mlp\_training.py}:
    \begin{itemize}
      \item \texttt{torch.manual\_seed(42)}
      \item Adam, lr\,=\,0.1, 500 épocas
      \item MSELoss final $\approx$ 0.0001
    \end{itemize}
    \vspace{0.5em}
    \textbf{Saídas}:
    \begin{enumerate}
      \item \texttt{mlp\_model.onnx} --- padrão industria
      \item \texttt{mlp\_model.pth} --- checkpoint
      \item \texttt{mlp\_weights.h} --- fallback float
    \end{enumerate}
    \vspace{0.3em}
    \begin{alertblock}{Por que ONNX?}
      Compatível com HuggingFace, ONNX Runtime
      e ferramentas de verificação formais.
    \end{alertblock}

    \column{0.5\textwidth}
    \begin{lstlisting}[style=pystyle, title=Arquitetura MLP (forward)]
# 2 entradas -> 4 neuronios -> 1 saida
def forward(self, x):
    x = self.hidden(x)  # Linear(2,4)
    x = self.relu(x)    # ReLU
    x = self.output(x)  # Linear(4,1)
    x = self.sigmoid(x) # [0,1]
    return x

# Treinamento
optimizer = optim.Adam(lr=0.1)
for epoch in range(500):
    loss = criterion(model(X), y)
    loss.backward()
    optimizer.step()
# Loss final: 0.0001
    \end{lstlisting}
  \end{columns}
\end{frame}

% ──────────────────────────────────────────────────────────────────────────────
\begin{frame}[fragile]{Etapa 2 --- Extração via ONNX}
  \begin{columns}
    \column{0.48\textwidth}
    \textbf{onnx\_mlp\_extractor.py}:
    \begin{itemize}
      \item Lê o grafo ONNX com \texttt{onnx}
      \item Localiza os dois nós \texttt{Gemm}
      \item Extrai tensores por nome
    \end{itemize}
    \vspace{0.3em}
    \textbf{Grafo identificado}:
    \begin{center}
      \scriptsize
      \begin{tabular}{lll}
        \toprule
        Nó & Op & Tensores \\
        \midrule
        1 & Gemm    & \texttt{hidden.weight [4,2]} \\
          &         & \texttt{hidden.bias [4]} \\
        2 & ReLU    & --- \\
        3 & Gemm    & \texttt{output.weight [1,4]} \\
          &         & \texttt{output.bias [1]} \\
        4 & Sigmoid & --- \\
        \bottomrule
      \end{tabular}
    \end{center}

    \column{0.5\textwidth}
    \begin{lstlisting}[style=pystyle, title=onnx\_mlp\_extractor.py]
def extract_mlp_weights(path):
  model = onnx.load(path)
  # nome -> tensor NumPy (initializers)
  tensors = {i.name: np.frombuffer(
      i.raw_data, dtype=np.float32
    ).reshape(i.dims)
    for i in model.graph.initializer}
  gemm = [n for n in model.graph.node
          if n.op_type == "Gemm"]
  return {
    "w_hidden": tensors[gemm[0].input[1]].tolist(),
    "b_hidden": tensors[gemm[0].input[2]].tolist(),
    "w_out":    tensors[gemm[1].input[1]][0].tolist(),
    "b_out":    float(tensors[gemm[1].input[2]][0]),
  }
    \end{lstlisting}
  \end{columns}
\end{frame}

% ──────────────────────────────────────────────────────────────────────────────
\begin{frame}{Etapa 3 --- Quantização Q8.8}
  \begin{block}{Esquema Q8.8 --- scale\,=\,256}
    \small
    $\hat{w} = \mathrm{round}(w \times 256)$;\quad
    $\hat{x} = x \times 256$;\quad
    $a \otimes b = \lfloor a \cdot b / 256 \rfloor$
  \end{block}

  \vspace{0.2em}
  \begin{columns}
    \column{0.52\textwidth}
    \begin{center}
      \tiny
      \begin{tabular}{@{}lrrrrl@{}}
        \toprule
        & $h_0$ & $h_1$ & $h_2$ & $h_3$ & viés \\
        \midrule
        $W_h$ & $[-183,-182]$ & $[-750,750]$ & $[-555,-555]$ & $[-153,952]$ & --- \\
        $b_h$ & 181 & 0 & 555 & 153 & \\
        $W_o$ & $-154$ & $1064$ & $-888$ & $-639$ & 1037 \\
        \bottomrule
      \end{tabular}
    \end{center}
    \vspace{0.2em}
    \textbf{Validação Python pré-ESBMC}:
    \begin{center}
      \scriptsize
      \begin{tabular}{cc|rr|c}
        \toprule
        $x_1$ & $x_2$ & float & quant & status \\
        \midrule
        0 & 0 & $-5.39$ & $-1380$ & \textcolor{okgreen}{OK} \\
        0 & 1 & $+5.45$ & $+1395$ & \textcolor{okgreen}{OK} \\
        1 & 0 & $+4.05$ & $+1037$ & \textcolor{okgreen}{OK} \\
        1 & 1 & $-5.23$ & $-1340$ & \textcolor{okgreen}{OK} \\
        \bottomrule
      \end{tabular}
    \end{center}

    \column{0.45\textwidth}
    \begin{exampleblock}{Invariante injetada}
      \small
      Limita busca SMT:\\[0.2em]
      $\texttt{pre} \in [-1280,\; 1280]$\\[0.2em]
      \scriptsize Derivado dos pesos.\\
      Via \texttt{\_\_ESBMC\_assume}.
    \end{exampleblock}
  \end{columns}
\end{frame}

% ──────────────────────────────────────────────────────────────────────────────
\begin{frame}[fragile]{Etapa 4 --- Harness de Verificação}
  \begin{columns}[T]
    \column{0.56\textwidth}
    \begin{lstlisting}[style=cstyle, basicstyle=\scriptsize\ttfamily, title=mlp\_forward (Q8.8)]
/* pesos: ver slide anterior */
static int mlp_forward(int x1, int x2) {
  int h[4], i;
  for (i = 0; i < 4; i++) {
    int pre =
      (x1*qw_hidden[i][0])/256
     +(x2*qw_hidden[i][1])/256
     + qb_hidden[i];
    __ESBMC_assume(
      pre>=-1280 && pre<=1280);
    h[i] = pre>0 ? pre : 0; /*ReLU*/
  }
  int score = qb_out;
  for (i=0;i<4;i++)
    score += (h[i]*qw_out[i])/256;
  return score;
}
    \end{lstlisting}

    \column{0.42\textwidth}
    \begin{lstlisting}[style=cstyle, basicstyle=\scriptsize\ttfamily, title=main (assertivas)]
int main(void) {
  __ESBMC_assert(
    mlp_forward(0,0)<=0,
    "P1: XOR(0,0)=false");
  __ESBMC_assert(
    mlp_forward(256,256)<=0,
    "P2: XOR(1,1)=false");
  __ESBMC_assert(
    mlp_forward(0,256)>0,
    "P3: XOR(0,1)=true");
  __ESBMC_assert(
    mlp_forward(256,0)>0,
    "P4: XOR(1,0)=true");
}
    \end{lstlisting}
  \end{columns}
\end{frame}

% ──────────────────────────────────────────────────────────────────────────────
\begin{frame}{Resultado --- Verificação Formal do XOR}
  \begin{center}
    \Large\textcolor{okgreen}{\textbf{VERIFICATION SUCCESSFUL}}\\[0.3em]
    \normalsize Solver: Boolector \quad|\quad Tempo solver: \textbf{0.001\,s} \quad|\quad Total: 0.105\,s
  \end{center}
  \vspace{0.5em}
  \begin{columns}
    \column{0.48\textwidth}
    \begin{exampleblock}{O que foi provado}
      Para pesos treinados (PyTorch, seed=42),\\
      quantizados em Q8.8:\\[0.4em]
      \begin{tabular}{cl}
        \textcolor{okgreen}{\checkmark} P1 & \texttt{mlp(0,\;0)} $\leq 0$ \\
        \textcolor{okgreen}{\checkmark} P2 & \texttt{mlp(1,\;1)} $\leq 0$ \\
        \textcolor{okgreen}{\checkmark} P3 & \texttt{mlp(0,\;1)} $> 0$ \\
        \textcolor{okgreen}{\checkmark} P4 & \texttt{mlp(1,\;0)} $> 0$ \\
      \end{tabular}
    \end{exampleblock}

    \column{0.48\textwidth}
    \begin{alertblock}{Teste vs.\ Verificação Formal}
      \begin{itemize}
        \item \textbf{Teste}: funciona \emph{nessas} entradas
        \item \textbf{Verificação}: \emph{matematicamente impossível} falhar com esses pesos
      \end{itemize}
    \end{alertblock}
    \vspace{0.3em}
    \begin{center}
      \scriptsize
      Fonte: \texttt{mlp\_model.onnx}\\
      Extrator: \texttt{onnx\_mlp\_extractor.py}\\
      Quantização: Q8.8, scale\,=\,256
    \end{center}
  \end{columns}
\end{frame}

% ══════════════════════════════════════════════════════════════════════════════
\section{Caso de Estudo 2 --- FFN de LLM}
% ══════════════════════════════════════════════════════════════════════════════

\begin{frame}{FFN de Transformers --- Contexto}
  \begin{columns}
    \column{0.52\textwidth}
    \textbf{Bloco FFN} em transformer:
    \[
      \mathbf{y} = W_2 \cdot \mathrm{GeLU}(W_1 \mathbf{x} + \mathbf{b}_1) + \mathbf{b}_2
    \]
    \begin{itemize}
      \item GPT-2:\quad $d_\text{model}=768$,\; $d_\text{ff}=3072$
      \item LLaMA-7B:\; $d_\text{model}=4096$,\; $d_\text{ff}=11008$
    \end{itemize}
    \vspace{0.4em}
    \textbf{Desafio}: verificar dimensões reais é intratável.\\
    \textbf{Solução}: \emph{slice} para dimensão tratável
    ($d_m=4,\; d_f=8$) e provar sobre pesos reais do modelo.

    \column{0.45\textwidth}
    \begin{block}{Propriedades verificadas}
      \begin{enumerate}
        \item Sem overflow em multiplicação de matrizes
        \item Pré-ativações dentro de intervalo analítico
        \item Saída dentro de intervalo derivado
        \item Sem acesso fora dos limites de arrays
      \end{enumerate}
    \end{block}
  \end{columns}
\end{frame}

% ──────────────────────────────────────────────────────────────────────────────
\begin{frame}[fragile]{Método QNNVerifier --- Injeção de Intervalo Abstrato}
  \begin{block}{Técnica central (ACASXU\_Reluplex\_fxp)}
    \small Substituir ativação por variável simbólica com limites analíticos:
    $h_j = \texttt{nondet\_int()}$;\quad
    $\texttt{\_\_ESBMC\_assume}(h_j \in [\ell_j, u_j])$
  \end{block}

  \begin{columns}
    \column{0.5\textwidth}
    \textbf{Otimizações descobertas}:
    \begin{enumerate}
      \item \textbf{int32\_t} em vez de int64\_t
      \item \textbf{Código morto eliminado}: inputs desacoplados
      \item \textbf{fxp32\_mult puro}: produto $\ll 2^{31}$
      \item \textbf{Boolector} em vez de Z3
    \end{enumerate}
    \begin{exampleblock}{}
      \centering \textbf{18--27$\times$ mais rápido}
    \end{exampleblock}

    \column{0.48\textwidth}
    \begin{lstlisting}[style=cstyle, basicstyle=\scriptsize\ttfamily, title=Harness QNN abstrato]
int32_t h[D_FF];
h[0] = nondet_int();
__ESBMC_assume(
  h[0] >= -9 && h[0] <= 10);
/* ... por neuronio j */

int32_t out0 = 0;
out0 += fxp32_mult(W2_0[0], h[0]);
/* ... */
__ESBMC_assert(
  out0 >= -12 && out0 <= 12,
  "P1: output[0] bounded");
    \end{lstlisting}
  \end{columns}
\end{frame}

% ──────────────────────────────────────────────────────────────────────────────
\begin{frame}{Resultados --- Verificação FFN (GPT-2 e LLaMA)}
  \begin{center}
    \textbf{Método 1} --- Ponto Fixo Q8.8 com proxy ReLU
    \vspace{0.2em}

    \small
    \begin{tabular}{llll}
      \toprule
      \textbf{Harness} & \textbf{Dimensões} & \textbf{Solver} & \textbf{Resultado} \\
      \midrule
      \texttt{mini\_ffn\_fxp\_test.c} & 2$\times$4$\times$2 & Z3 & \textcolor{okgreen}{0.25\,s} \\
      \texttt{gpt2\_4x8.c}            & 4$\times$8$\times$4 & Z3 & \textcolor{okgreen}{12.4\,s} \\
      \texttt{llama\_4x8.c}           & 4$\times$8$\times$4 & Z3 & \textcolor{okgreen}{8.6\,s} \\
      \texttt{gpt2\_4x16.c}           & 4$\times$16$\times$4 & Z3 & \textcolor{failred}{timeout} \\
      \bottomrule
    \end{tabular}
  \end{center}

  \vspace{0.6em}
  \begin{center}
    \textbf{Método 2} --- QNNVerifier (ativação exata via intervalo abstrato)
    \vspace{0.2em}

    \small
    \begin{tabular}{lllll}
      \toprule
      \textbf{Harness} & \textbf{Dims} & \textbf{Ativação} & \textbf{Antes} & \textbf{Depois} \\
      \midrule
      \texttt{gpt2\_2x4\_qnn.c}     & 2$\times$4$\times$2  & GeLU & ---        & \textcolor{okgreen}{0.30\,s} \\
      \texttt{gpt2\_4x8\_qnn.c}     & 4$\times$8$\times$4  & GeLU & 94--115\,s & \textcolor{okgreen}{5.3\,s} \\
      \texttt{llama-7b\_4x8\_qnn.c} & 4$\times$8$\times$4  & SiLU & 115\,s     & \textcolor{okgreen}{4.3\,s} \\
      \texttt{gpt2\_4x16\_qnn.c}    & 4$\times$16$\times$4 & GeLU & timeout    & \textcolor{okgreen}{27\,s} \\
      \bottomrule
    \end{tabular}
  \end{center}
  \vspace{0.2em}
  \begin{center}
    \scriptsize ESBMC 6.8.0, Boolector 3.2, Linux x86\_64
  \end{center}
\end{frame}

% ══════════════════════════════════════════════════════════════════════════════
\section{Contribuições Desta Sessão}
% ══════════════════════════════════════════════════════════════════════════════

\begin{frame}{Contribuições Desenvolvidas Nesta Sessão}
  \begin{enumerate}
    \item \textbf{Treinamento real com PyTorch}\\
          Pesos agora são treinados do zero (seed=42, 500 épocas, loss $\approx 0.0001$)

    \vspace{0.3em}
    \item \textbf{Export ONNX funcionando} (\texttt{onnxscript} instalado)\\
          \texttt{mlp\_training.py} exporta \texttt{mlp\_model.onnx} automaticamente

    \vspace{0.3em}
    \item \textbf{onnx\_mlp\_extractor.py} --- módulo novo\\
          Extrai pesos do grafo ONNX inspecionando nós \texttt{Gemm}

    \vspace{0.3em}
    \item \textbf{quantize\_mlp.py refatorado}\\
          \texttt{-{}-source onnx} (padrão) lê ONNX; \texttt{-{}-source h} usa fallback\\
          Valida 4 casos XOR em Python antes de gerar o harness

    \vspace{0.3em}
    \item \textbf{Pipeline end-to-end completo e automatizado}\\
          3 comandos: treinar $\to$ quantizar $\to$ verificar

    \vspace{0.3em}
    \item \textbf{Otimizações QNNVerifier (FFN/LLM)}\\
          int32 + dead-code removal + fxp32\_mult + Boolector = \textbf{18--27$\times$ speedup}
  \end{enumerate}
\end{frame}

% ──────────────────────────────────────────────────────────────────────────────
\begin{frame}[fragile]{Pipeline Final em 3 Comandos}
  \begin{center}
  \begin{tikzpicture}
    \node[pipeblue]    (tr) at (0,0)   {\texttt{python mlp\_training.py}};
    \node[pipeorange]  (on) at (5,0)   {\texttt{mlp\_model.onnx}};
    \node[pipeyellow]  (qp) at (10,0)  {\texttt{python quantize\_mlp.py}};
    \node[pipered]     (hc) at (10,-2) {\texttt{verify\_mlp\_qnn.c}};
    \node[pipegreen]   (es) at (5,-2)  {\texttt{esbmc ... -{}-boolector}};
    \node[pipegreendk] (ok) at (0,-2)  {\textcolor{okgreen}{\large\textbf{SUCCESSFUL}}};

    \draw[piparr] (tr) -- (on)  node[above,piplbl,midway]{treina + exporta};
    \draw[piparr] (on) -- (qp)  node[above,piplbl,midway]{lê via extrator};
    \draw[piparr] (qp) -- (hc)  node[right,piplbl,midway]{gera C};
    \draw[piparr] (hc) -- (es);
    \draw[piparr] (es) -- (ok)  node[below,piplbl,midway]{0.001\,s solver};
  \end{tikzpicture}
  \end{center}

  \vspace{0.5em}
  \begin{columns}
    \column{0.48\textwidth}
    \begin{exampleblock}{Propriedades provadas}
      \begin{tabular}{cl}
        \textcolor{okgreen}{\checkmark} P1 & XOR(0,0) = false \\
        \textcolor{okgreen}{\checkmark} P2 & XOR(1,1) = false \\
        \textcolor{okgreen}{\checkmark} P3 & XOR(0,1) = true \\
        \textcolor{okgreen}{\checkmark} P4 & XOR(1,0) = true \\
      \end{tabular}
    \end{exampleblock}

    \column{0.48\textwidth}
    \begin{block}{Arquivos principais}
      \scriptsize
      \texttt{mlp\_training.py}\\
      \texttt{onnx\_mlp\_extractor.py} \textbf{(novo)}\\
      \texttt{quantize\_mlp.py} \textbf{(refatorado)}\\
      \texttt{verify\_mlp\_qnn.c} \textbf{(gerado)}
    \end{block}
  \end{columns}
\end{frame}

% ══════════════════════════════════════════════════════════════════════════════
\section{Próximos Passos}
% ══════════════════════════════════════════════════════════════════════════════

\begin{frame}[fragile]{Verificação de Neurônios Mortos --- Resultados}
  \begin{columns}
    \column{0.48\textwidth}
    \textbf{Método}: para cada $h_i$, entradas \textbf{simbólicas} em $[0,256]^2$:
    \begin{lstlisting}[style=cstyle]
int x1 = nondet_int();
int x2 = nondet_int();
__ESBMC_assume(x1>=0 && x1<=256);
__ESBMC_assume(x2>=0 && x2<=256);
/* FAILED  = neuronio VIVO
   SUCCESS = neuronio MORTO */
__ESBMC_assert(h[i] == 0,
  "neuronio i sempre morto?");
    \end{lstlisting}

    \vspace{0.15em}
    \begin{exampleblock}{Modelo real (seed=42)}
      \scriptsize
      \begin{tabular}{@{}crl@{}}
        \toprule
        $h_i$ & $b$ & ESBMC \\
        \midrule
        0 & $+181$ & \textcolor{okgreen}{\textbf{VIVO}} \\
        1 & $0$    & \textcolor{okgreen}{\textbf{VIVO}} \\
        2 & $+555$ & \textcolor{okgreen}{\textbf{VIVO}} \\
        3 & $+153$ & \textcolor{okgreen}{\textbf{VIVO}} \\
        \bottomrule
      \end{tabular}
      \quad Nenhum morto.
    \end{exampleblock}

    \column{0.5\textwidth}
    \begin{alertblock}{Demonstração: neurônio morto artificial}
      \small
      Modifica $b[1]: 0 \to -1000$\\[0.1em]
      \textbf{Prova}: $\max(\text{pre}_1)=750-1000=-250\leq 0$\\
      $\Rightarrow h_1=0$ para toda entrada
    \end{alertblock}
    \vspace{0.1em}
    \begin{exampleblock}{Resultado com $b[1]=-1000$}
      \scriptsize
      \begin{tabular}{@{}crl@{}}
        \toprule
        $h_i$ & $b$ & ESBMC \\
        \midrule
        0 & $+181$ & \textcolor{okgreen}{\textbf{VIVO}} \\
        1 & $-1000$ & \textcolor{failred}{\textbf{MORTO}} \\
        2 & $+555$ & \textcolor{okgreen}{\textbf{VIVO}} \\
        3 & $+153$ & \textcolor{okgreen}{\textbf{VIVO}} \\
        \bottomrule
      \end{tabular}
      \quad \textcolor{failred}{VERIFICATION SUCCESSFUL}
    \end{exampleblock}
  \end{columns}
\end{frame}

% ──────────────────────────────────────────────────────────────────────────────
\begin{frame}{Roadmap e Conclusão}
  \begin{columns}
    \column{0.48\textwidth}
    \textbf{Roadmap}:
    \begin{itemize}
      \item \textcolor{okgreen}{\checkmark} MLP XOR treinado (PyTorch real)
      \item \textcolor{okgreen}{\checkmark} Pipeline ONNX end-to-end
      \item \textcolor{okgreen}{\checkmark} FFN GPT-2 / LLaMA verificado
      \item \textcolor{okgreen}{\checkmark} QNNVerifier 18--27$\times$ mais rápido
      \item \textcolor{okgreen}{\checkmark} Verificação de neurônios mortos
      \item \textcolor{gray}{$\square$} Propriedades de robustez ($\ell_\infty$)
      \item \textcolor{gray}{$\square$} Frama-C EVA para bounds mais justos
      \item \textcolor{gray}{$\square$} Escalar para $d_\text{ff} > 16$
    \end{itemize}

    \column{0.48\textwidth}
    \begin{exampleblock}{Impacto}
      \begin{itemize}
        \item Modelo verificado = \textbf{pronto para deploy seguro}
        \item Pipeline reutilizável para qualquer MLP ONNX
        \item Base para escalar a LLMs maiores (slices)
        \item Speedup de 18--27$\times$ vs.\ estado anterior
      \end{itemize}
    \end{exampleblock}

    \vspace{0.5em}
    \begin{center}
      \Large\textcolor{esbmcblue}{\texttt{VERIFICATION SUCCESSFUL}}
    \end{center}
  \end{columns}
\end{frame}

% ──────────────────────────────────────────────────────────────────────────────
\begin{frame}[plain]
  \begin{center}
    \vspace{2cm}
    {\Huge Obrigado!}\\[1.5em]
    \large
    Código disponível em:\\[0.5em]
    \texttt{pibic/teste\_mlp/}\\[0.2em]
    \texttt{pibic/cases/llm\_ffn\_verification/}\\[1.5em]
    \normalsize
    Branch: \texttt{claude/analyze-pibic-code-KfVX4}
  \end{center}
\end{frame}

\end{document}
